\documentclass{beamer}
\usepackage[english]{babel}
\usepackage{graphicx}
\usepackage{booktabs}


\title{Weekly report}
\author{Natalia Trigo Tomasevich}
\date{\today}

\begin{document}


\frame{\titlepage}

\begin{frame}{From last week}
    \begin{itemize}
        \item Try to prove that the solution was a threshold policy
        \item The fluid model
    \end{itemize}
\end{frame}

\begin{frame}{The base case}
    \centering
    \includegraphics[width=\textwidth,height=0.82\textheight,keepaspectratio]{../../forgetting/outputs/advantage_by_p0_and_rho/p0_0.5/figures/advantage_by_rho_overlay.pdf}
\end{frame}

\begin{frame}{No unimodality}
    \centering
    \includegraphics[width=\textwidth,height=0.82\textheight,keepaspectratio]{../../forgetting/outputs/numerical_validations/discrete_gap_counterexample/counterexample.png}
\end{frame}

\begin{frame}{A weird case}
    \centering
    \includegraphics[width=\textwidth,height=0.82\textheight,keepaspectratio]{../../forgetting/outputs/threshold_parameter_map/figures/05_a_rho_one_lower_case.pdf}
\end{frame}

\begin{frame}
    \centering
    \includegraphics[width=\textwidth,height=0.94\textheight,keepaspectratio]{../../forgetting/outputs/threshold_parameter_map/figures/01_policy_phase_maps.pdf}
\end{frame}

\begin{frame}{About the thresholds}
    \begin{itemize}
        \item We don't always have thresholds. It depends on the parameters of the problem 
        \item For some parameters we do have thresholds
        \item We could have up to 2 threholds. 
        \item They are always the same patterns $$1\rightarrow2 , \qquad 1\rightarrow2\rightarrow1, \qquad 2\rightarrow1$$
    \end{itemize}
\end{frame}

\begin{frame}{We don't always have thresholds}
    If we have $$\mathcal A^1_\rho =\max_{z\in\mathcal K_\rho}
    \left\{-\Delta c+\gamma\Delta p \left[U_\rho^1(T_\rho^s z)-U_\rho^1(T_\rho^f z)\right]\right\} <0$$
    Then for all $z$ we have that to use product 1 is optimal.  \\
    If we have   $$\mathcal A^2_\rho =\min_{z\in\mathcal K_\rho}
    \left\{-\Delta c+\gamma\Delta p \left[U_\rho^2(T_\rho^s z)-U_\rho^2(T_\rho^f z)\right]\right\} >0$$ Then for all $z$ we have that to use product 2 is optimal. 
\end{frame}

\begin{frame}{We have thresholds}
    We have $$\mathcal A^1_\rho =\max_{z\in\mathcal K_\rho}
    \left\{-\Delta c+\gamma\Delta p \left[U_\rho^1(T_\rho^s z)-U_\rho^1(T_\rho^f z)\right]\right\} \mathrel{\textcolor{red}{>}}0$$
    $$\iff$$
    $$\mathcal A_\rho = \max_{z\in\mathcal K_\rho}
    \underbrace{\left\{-\Delta c+\gamma\Delta p \left[V_\rho(T_\rho^s z)-V_\rho(T_\rho^f z)\right]\right\}}_{a_\rho(z)} >0$$
\end{frame}

\begin{frame}{We have thresholds}
    I proved that
    $$\mathcal A^2_\rho =\min_{z\in\mathcal K_\rho}
    \left\{-\Delta c+\gamma\Delta p \left[U_\rho^2(T_\rho^s z)-U_\rho^2(T_\rho^f z)\right]\right\} \mathrel{\textcolor{red}{<}}0$$
    $$\iff$$
    $$\mathcal A'_\rho = \min_{z\in\mathcal K_\rho}
    \left\{-\Delta c+\gamma\Delta p \left[V_\rho(T_\rho^s z)-V_\rho(T_\rho^f z)\right]\right\} <0$$ 
\end{frame}

\begin{frame}{Threshold area}
    If $\mathcal A^1_\rho>0$ and $\mathcal A^2_\rho<0$ we have at least one threshold,because that condition means that $\exists z_1,z_2\in \mathcal K_\rho$ such that $a_\rho(z_2)<0<a_\rho(z_1)$. Which means, by continuity and intermediate value that $\exists z^*\in (min\{z_1,z_2\},max\{z_1,z_2\})$ such that $a_\rho(z^*)=0.$
\end{frame}

\begin{frame}{Conditions over $c_2$}
    \begin{itemize}
        \item If we fix all the parameters except for $c_2$, with the conditions $\mathcal A^1_\rho>0$ and $\mathcal A^2_\rho<0$ we can find an upper bound and lower bound for $c_2$ where we can find thresholds. 
        \item The proof is not revised yet.  
    \end{itemize}
\end{frame}

\begin{frame}{The fluid model}
    Proceed to progress note...
\end{frame}

\end{document}